\documentclass[10pt,aspectratio=169]{beamer}

\usetheme{Warsaw}
% Thème couleur personnalisé marron/doré
\setbeamertemplate{navigation symbols}{}
\setbeamertemplate{headline}{%
  \begin{beamercolorbox}[wd=\paperwidth,ht=2.8ex,dp=1ex,center]{section in head/foot}%
    \scriptsize%
    \insertsectionnavigationhorizontal{\paperwidth}{}{}%
  \end{beamercolorbox}%
}
\setbeamercolor{section in head/foot}{bg=marron, fg=white}
\setbeamertemplate{footline}{%
  \leavevmode%
  \hbox{%
    \begin{beamercolorbox}[wd=0.18\paperwidth,ht=2.6ex,dp=1ex,center]{author in head/foot}%
      \scriptsize \insertframenumber/\inserttotalframenumber
    \end{beamercolorbox}%
    \begin{beamercolorbox}[wd=0.50\paperwidth,ht=2.6ex,dp=1ex,center]{author in head/foot}%
      \scriptsize Sié Rachid TRAORÉ
    \end{beamercolorbox}%
    \begin{beamercolorbox}[wd=0.32\paperwidth,ht=2.6ex,dp=1ex,center]{title in head/foot}%
      \scriptsize Élève Ingénieur Statisticien Économiste
    \end{beamercolorbox}%
  }%
}
\setbeamertemplate{caption}[numbered]
\setbeamertemplate{blocks}[default]

% Couleurs personnalisées
\definecolor{marron}{RGB}{101, 83, 45}
\definecolor{dore}{RGB}{184, 134, 11}
\setbeamercolor{frametitle}{bg=marron, fg=white}
\setbeamercolor{title}{bg=marron, fg=white}
\setbeamercolor{subtitle}{fg=white}
\setbeamercolor{structure}{fg=marron}
\setbeamercolor{palette primary}{bg=marron, fg=white}
\setbeamercolor{palette secondary}{bg=marron!85!black, fg=white}
\setbeamercolor{palette tertiary}{bg=dore!85!black, fg=white}
\setbeamercolor{palette quaternary}{bg=marron!70!black, fg=white}
\setbeamercolor{alerted text}{fg=dore!85!black}
\setbeamercolor{block title}{bg=marron, fg=white}
\setbeamercolor{block body}{bg=marron!6, fg=black}
\setbeamercolor{alertblock title}{bg=dore!85!black, fg=white}
\setbeamercolor{alertblock body}{bg=dore!10, fg=black}

\usepackage[T1]{fontenc}
\usepackage[utf8]{inputenc}
\usepackage[french]{babel}
\usepackage{amsmath,amssymb}
\usepackage{booktabs}
\usepackage{graphicx}
\usepackage{tikz}
\usetikzlibrary{positioning}

\title[Transferts de migrants et pauvreté des enfants]{%
  Impact des transferts de migrants sur la pauvreté\\
  multidimensionnelle des enfants au Sénégal}
\subtitle{}
\author[]{Sié Rachid TRAORÉ}
\institute[]{}
\date{}


\begin{document}

\begingroup
\setbeamertemplate{headline}{}
\begin{frame}
  \centering
  \vspace{0.45cm}
  \begin{beamercolorbox}[sep=8pt,center,rounded=true,shadow=true]{title}
    {\LARGE\bfseries\color{white}
      Impact des transferts de migrants sur la pauvreté
      multidimensionnelle des enfants au Sénégal\par}
  \end{beamercolorbox}
  \vspace{0.55cm}
  {\large Sié Rachid TRAORÉ\par}\vspace{0.15cm}
  {\small Élève Ingénieur Statisticien Économiste\par}
  \vfill
  {\small Sous la direction de Prénom NOM\par}\vspace{0.12cm}
  {\small Titre --- Institution\par}
  \vspace{0.25cm}
  {\small Juin 2026\par}
\end{frame}
\endgroup

\begin{frame}{Plan de la présentation}
  \tableofcontents
\end{frame}


\section[Intro]{Introduction}

\begin{frame}{Contexte et motivation}
  \begin{block}{Place des transferts de migrants}
    \begin{itemize}
      \item Les transferts de fonds représentent environ \textbf{10--12\,\%} du PIB sénégalais.
      \item Première source de devises, supérieure à l'APD et aux IDE.
      \item Plus de 3 millions de Sénégalais vivent à l'étranger.
    \end{itemize}
  \end{block}
  \begin{alertblock}{Pauvreté multidimensionnelle des enfants}
    \begin{itemize}
      \item Près de 50\,\% des enfants souffrent de privations multiples (ANSD, 2022).
      \item La pauvreté monétaire ne suffit pas à capturer toutes ces dimensions.
    \end{itemize}
  \end{alertblock}
\end{frame}

\begin{frame}{Problématique et objectifs}
  \begin{block}{Question centrale}
    Dans quelle mesure les transferts de migrants réduisent-ils
    la \textbf{pauvreté multidimensionnelle} des enfants au Sénégal ?
  \end{block}
  \pause
  \begin{block}{Objectif général}
    Évaluer l'impact causal des transferts de migrants sur la pauvreté
    multidimensionnelle des enfants à partir des données EHCVM I \& II.
  \end{block}
  \pause
  \begin{block}{Hypothèses}
    \begin{enumerate}
      \item[H1] Les transferts réduisent significativement la pauvreté multidimensionnelle des enfants.
      \item[H2] L'effet est hétérogène selon les dimensions et le milieu de résidence.
      \item[H3] Le PSM-DD contrôle efficacement les biais de sélection.
    \end{enumerate}
  \end{block}
\end{frame}

\section[Revue]{Revue de la littérature}

\begin{frame}{Transferts de migrants et bien-être : théories}
  \begin{itemize}
    \item \textbf{Altruisme} (Lucas \& Stark, 1985) : solidarité familiale comme moteur principal.
    \item \textbf{Contrat familial implicite} : les transferts compensent l'investissement consenti dans la migration.
    \item \textbf{Diversification des risques} : le ménage envoie un membre migrer pour sécuriser ses revenus.
  \end{itemize}
  \begin{block}{Effets empiriques documentés}
    \begin{itemize}
      \item Réduction de la pauvreté monétaire (Adams \& Page, 2005).
      \item Amélioration de la scolarisation et de la santé des enfants (Yang, 2008).
      \item Effets plus marqués en milieu rural et pour les ménages les plus pauvres.
    \end{itemize}
  \end{block}
\end{frame}

\begin{frame}{Pauvreté multidimensionnelle : approche Alkire-Foster}
  \begin{itemize}
    \item Le revenu ou la consommation ne capturent pas toutes les privations des enfants.
    \item L'approche Alkire-Foster (2011) identifie les privations dans plusieurs dimensions.
    \item Double seuil : intra-dimensionnel (seuil de privation) et inter-dimensionnel ($k$).
    \item Mesure $M_0 = H \times A$ : incidence $\times$ intensité.
  \end{itemize}
  \begin{alertblock}{Lacune identifiée}
    Peu d'études combinent PSM et double différence pour identifier l'effet causal
    des transferts sur la pauvreté \textbf{multidimensionnelle} des enfants au Sénégal.
  \end{alertblock}
\end{frame}

\section[Méthodo]{Méthodologie}

\begin{frame}{Construction de l'IPM-Enfant}
  \begin{block}{Dimensions et indicateurs (poids égaux : $w_j = 1/6$)}
    \centering\smallskip
    \begin{tabular}{lll}
      \toprule
      Dimension     & Indicateur              & Seuil de privation \\
      \midrule
      Éducation     & Fréquentation scolaire  & Non scolarisé (6--17 ans) \\
                    & Retard scolaire         & Retard $\geq$ 2 ans \\
      Santé         & Suivi nutritionnel      & Aucun suivi / 12 mois \\
                    & Vaccination             & Vaccination incomplète \\
      Niveau de vie & Eau potable             & Eau non améliorée \\
                    & Assainissement          & Toilettes non améliorées \\
      \bottomrule
    \end{tabular}
  \end{block}
  \begin{equation*}
    M_0 = H \times A, \quad k = \tfrac{1}{3}
  \end{equation*}
\end{frame}

\begin{frame}{Stratégie d'identification : PSM-Double différence}
  \begin{block}{Problème : biais de sélection}
    Les ménages bénéficiaires de transferts diffèrent systématiquement des non-bénéficiaires.
  \end{block}
  \begin{block}{Étape 1 --- Score de propension (probit)}
    \begin{equation*}
      p(X_i) = P(D_i = 1 \mid X_i) = \Phi(X_i'\beta)
    \end{equation*}
    Appariement $k$-plus proches voisins ($k=4$), vérification du support commun.
  \end{block}
  \begin{alertblock}{Étape 2 --- Double différence après appariement (PSM-DD)}
    \begin{equation*}
      \hat{\tau}_{PSM\text{-}DD}
      = \frac{1}{n_T}\sum_{i\in\mathcal{T}}
        \!\left[\Delta Y_i - \sum_{j\in\mathcal{C}} w_{ij}\,\Delta Y_j\right]
    \end{equation*}
    Contrôle simultané des biais observables et des effets fixes inobservés.
  \end{alertblock}
\end{frame}

\begin{frame}{Données : EHCVM I et II}
  \begin{itemize}
    \item \textbf{EHCVM I (2018--2019)} : première vague harmonisée UEMOA au Sénégal.
    \item \textbf{EHCVM II (2021--2022)} : deuxième vague constituant un panel.
    \item \textbf{Variable de traitement} : réception de transferts de migrants dans les 12 derniers mois.
    \item \textbf{Population cible} : enfants de 0 à 17 ans dans les ménages appariés.
    \item \textbf{Covariables} : âge, sexe et éducation du chef de ménage ; taille du ménage ; milieu de résidence ; région.
  \end{itemize}
\end{frame}

\section[Résultats]{Résultats empiriques}

\begin{frame}{Pauvreté multidimensionnelle des enfants : état des lieux}
  \begin{block}{IPM-Enfant par vague (à compléter)}
    \begin{itemize}
      \item IPM-Enfant en 2018--2019 : \dots
      \item IPM-Enfant en 2021--2022 : \dots
      \item Dimension la plus répandue : \dots
    \end{itemize}
  \end{block}
  \begin{block}{Qualité de l'appariement}
    \begin{itemize}
      \item Support commun : pas de ménages hors zone de recouvrement.
      \item Différences standardisées $< 5\,\%$ après appariement.
    \end{itemize}
  \end{block}
\end{frame}

\begin{frame}{Impact des transferts sur l'IPM-Enfant}
  \begin{table}
    \centering
    \small
    \begin{tabular}{lcc}
      \toprule
      Méthode & ATT & $p$-valeur \\
      \midrule
      PSM -- $k$-ppv ($k=1$)  & \dots & \dots \\
      PSM -- $k$-ppv ($k=4$)  & \dots & \dots \\
      PSM -- noyau             & \dots & \dots \\
      \midrule
      \textbf{PSM-DD}          & \dots & \dots \\
      \bottomrule
    \end{tabular}
  \end{table}
  \begin{itemize}
    \item Un ATT négatif indique que les transferts \textbf{réduisent} l'IPM-Enfant.
    \item Hétérogénéité : effets plus marqués en milieu \dots
  \end{itemize}
\end{frame}

\begin{frame}{Validation des hypothèses}
  \begin{table}
    \centering
    \begin{tabular}{clc}
      \toprule
      Hypothèse & Résultat & Décision \\
      \midrule
      H1 & Réduction significative de l'IPM-Enfant & \dots \\
      H2 & Hétérogénéité selon milieu et dimension  & \dots \\
      H3 & PSM-DD contrôle les biais                & \dots \\
      \bottomrule
    \end{tabular}
  \end{table}
\end{frame}


\section[Conclusion]{Conclusion et recommandations}

\begin{frame}{Conclusion et recommandations}
  \begin{columns}[T]
    \begin{column}{0.48\textwidth}
      \begin{block}{Conclusion}
        \begin{itemize}
          \item Les transferts de migrants \dots la pauvreté multidimensionnelle des enfants.
          \item Effets hétérogènes selon le milieu et la dimension.
          \item Résultats robustes aux choix méthodologiques.
        \end{itemize}
      \end{block}
    \end{column}
    \begin{column}{0.48\textwidth}
      \begin{block}{Recommandations}
        \begin{itemize}
          \item Réduire les coûts de transaction des transferts formels.
          \item Orienter les transferts vers l'éducation et la santé des enfants.
          \item Cibler les non-bénéficiaires dans les programmes sociaux.
          \item Renforcer les statistiques sur les migrations.
        \end{itemize}
      \end{block}
    \end{column}
  \end{columns}
\end{frame}

\begin{frame}{Limites et pistes de recherche}
  \begin{columns}[T]
    \begin{column}{0.48\textwidth}
      \begin{alertblock}{Limites}
        \begin{itemize}
          \item Panel non cylindré entre les deux vagues.
          \item Hypothèse de tendances parallèles invérifiable sur l'inobservable.
          \item Mesure imparfaite des montants transférés.
        \end{itemize}
      \end{alertblock}
    \end{column}
    \begin{column}{0.48\textwidth}
      \begin{block}{Pistes}
        \begin{itemize}
          \item Analyser les mécanismes de transmission (dépenses en éducation, santé).
          \item Étendre à d'autres pays UEMOA couverts par l'EHCVM.
          \item Intégrer la dimension genre de l'enfant.
        \end{itemize}
      \end{block}
    \end{column}
  \end{columns}
\end{frame}


\begin{frame}
  \begin{center}
    \Huge Merci de votre attention
    \vspace{0.6cm}

    \large Questions et échanges
  \end{center}
\end{frame}

\end{document}
